\documentclass[12pt]{article}

%% Basic document formatting
\usepackage{amsmath, amsthm, amssymb, amsfonts}
\usepackage{mathtools}
\usepackage{xspace}
\usepackage{thmtools}
\usepackage{graphicx}
\usepackage{setspace}
\usepackage{fancyhdr}
\usepackage{titling}
\usepackage[left=0.4in,right=0.4in,top=1in,bottom=1in]{geometry}
\usepackage{float}
\usepackage{hyperref}
\usepackage{mathrsfs}
\usepackage{tabularx}
\usepackage[utf8]{inputenc}
\usepackage[english]{babel}
\usepackage{framed}
\usepackage[dvipsnames]{xcolor}
\usepackage{environ}
\usepackage{tcolorbox}
\tcbuselibrary{theorems,skins,breakable}

\usepackage{enumitem}
\setlist[enumerate]{leftmargin=*}
\setlist[enumerate,1]{labelindent=\parindent}
\setlist[enumerate,2]{labelindent=0pt}

% Blackboard Bold
\newcommand{\Z}{\mathbb{Z}}                 % integers
\newcommand{\Zpl}{\mathbb{Z}_{+}}           % positive integers
\newcommand{\Znn}{\mathbb{Z}_{\geq 0}}      % nonnegative integers. 
\newcommand{\Q}{\mathbb{Q}}                 % rationals
\newcommand{\Qpl}{\mathbb{Q}_{+}}           % positive Rationals
\newcommand{\R}{\mathbb{R}}                 % reals
\newcommand{\Rpl}{\mathbb{R}_{+}}           % positive Reals
\newcommand{\C}{\mathbb{C}}                 % complex numbers
\newcommand{\F}{\mathbb{F}}                 % field

% Words
\newcommand{\st}{\text{ such that }}
\newcommand{\wrt}{\text{ with respect to }}
\newcommand{\with}{\text{ with }}
\newcommand{\ie}{\text{, i.e. }}
\newcommand*{\ditto}{\text{---}\texttt{"}\text{---}}

% Operators
\newcommand{\abs}[1]{\left|#1\right|}                   % absolute value
\newcommand{\norm}[1]{\abs{\abs{#1}}}
\newcommand{\floor}[1]{\left\lfloor #1 \right\rfloor}   % floor
\newcommand{\ceil}[1]{\left\lceil #1 \right\rceil}      % ceiling
\newcommand{\powset}[1]{\mathscr{P}(#1)}

% Category Theory 
\renewcommand{\hom}{\text{Hom}}
\newcommand{\homspace}[3]{\hom_{#1}(#2, #3)}
\newcommand{\coproduct}[9]{
  \begin{tikzcd}[ampersand replacement=\&]
      \& \& #1 \arrow[dll, bend right, "#6"] \arrow[dl, "#5"] \\
   #4 \& #3 \arrow[l, "#9"] \\
      \& \& #2 \arrow[ull, bend left, "#8"] \arrow[ul, "#7"]
  \end{tikzcd}
}

\newcommand{\product}[9]{ 
  \begin{tikzcd}[ampersand replacement=\&]
      \& \& #3 \\
      #1 \arrow[r, "#7"] \arrow[urr, bend left, "#5"] \arrow[drr, bend right, "#6"] \& #2 \arrow[ur, "#8"] \arrow[dr, "#9"] \\ 
      \& \& #4
  \end{tikzcd}
}


% Algebra 
\newcommand{\Syl}[2]{\text{Syl}_{#1}{(#2)}}           % set of Sylow-p subgroups 
\newcommand{\<}{\langle}                % \<x\>, subgroup generated by x 
\renewcommand{\>}{\rangle}
\renewcommand{\index}[2]{[#1 : #2]}
\newcommand{\id}{\text{id}}             % identity element
\newcommand{\lhdneq}{%
  \mathrel{\ooalign{$\lneq$\cr\raise.22ex\hbox{$\lhd$}\cr}}}
\newcommand{\order}[1]{\text{o}(#1)}    % order of an element
\newcommand{\spec}{\operatorname{Spec}}
\newcommand{\rad}{\operatorname{rad}}   % radical of an ideal 
\newcommand{\sym}[1]{\mathfrak{S}_{#1}}
\newcommand{\aut}[1]{\text{Aut}(#1)} 
\newcommand{\gal}[2]{\text{Gal}(#1 / #2)} 
\newcommand{\inn}[1]{\text{Inn}(#1)} 
\let\oldcong\cong
\let\oldequiv\equiv
\renewcommand{\cong}{\oldequiv}
\renewcommand{\equiv}{\oldcong}

\newcommand{\triangleleftneq}{\mathrel{\ooalign{%
  $\trianglelefteq$\cr
  \hidewidth\raise0.06ex\hbox{$\not\mkern4mu$}\hidewidth\cr
}}}

\makeatletter % cycle
\newcommand{\cyc}[1]{(\mathbf{\cyc@process#1\relax})}
\def\cyc@process#1#2\relax{%
	#1%
	\ifx\relax#2\relax
	\else
		\,\cyc@process#2\relax
	\fi
}
\makeatother

\newcommand{\GL}[2]{\text{GL}_{#1}(#2)}                             % general linear group
\newcommand{\SL}[2]{\text{SL}_{#1}(#2)}                             % special linear group
\newcommand{\gl}[2]{\mathfrak{gl}_{#1}(#2)}
\renewcommand{\sl}[2]{\mathfrak{sl}_{#1}(#2)}
\newcommand{\so}[2]{\mathfrak{so}_{#1}(#2)}
\newcommand{\su}[1]{\mathfrak{su}(#1)}

% Linear Algebra
\newcommand{\bmat}[1]{\begin{bmatrix}#1\end{bmatrix}}   % bracketed matrix
\newcommand{\rank}{\operatorname{rank}}                 % rank
\newcommand{\nullity}{\operatorname{nullity}}           % nullity
\newcommand{\tr}{\operatorname{tr}}

% Combinatorics
\newcommand{\qnum}[1]{\left[#1\right]_q}                % q-analog
\newcommand{\qfact}[1]{\left[#1\right]!_q}              % q-analog factorial 
\newcommand{\qbinom}[2]{\binom{#1}{#2}_q}               % Gaussian binomial coefficient

% Topology/Analysis
\newcommand{\ball}[2]{\text{B}_{#1}(#2)}  % B_r(x): open r-balls around x
\newcommand{\diam}{\text{diam}}           % diamter of a set in metric space 
\newcommand{\lowint}[2]{\underline{\int_{#1}^{#2}}}
\newcommand{\upint}[2]{\overline{\int}_{#1}^{#2}}

% Misc Notation
\newcommand{\oneton}[1]{\{1, \ldots, #1\}}
\newcommand{\defeq}{\vcentcolon=}   % :=
\newcommand{\eqdef}{=\vcentcolon}   % =: 
\renewcommand{\bf}[1]{\textbf{#1}}
\newcommand{\im}{\operatorname{im}}
\newcommand{\fnrest}[2]{\left.#1\right|_{#2}}
\newcommand{\isom}{\overset{\sim}{\rightarrow}} 


\renewcommand{\thesection}{\arabic{section}.}
\renewcommand{\thesubsection}{(\alph{subsection})}

\newtheorem{claim}{Claim}
\newtheorem{theorem}{Theorem}
\newtheorem{proposition}{Proposition}
\newtheorem*{lemma}{Lemma}

%% Headers & title setup
\newcommand{\course}{Math140C}
\newcommand{\myname}{Jay Ser}
\setlength{\headheight}{14.5pt}
\pagestyle{fancy}
\fancyhf{}
\renewcommand{\headrulewidth}{0.4pt}
\lhead{\course}
\rhead{\myname}
\cfoot{\thepage}
\setlength{\droptitle}{-4em} 
\title{\course\ - HW \#7}
\author{\myname}
\date{2026.06.07} 

\begin{document}
\maketitle
\thispagestyle{fancy}

%------------------------------------------------------------------------------%

\newcommand{\piInt}{[-\pi, \pi]} 
\newcommand{\Ltwospace}[2]{\mathscr{L}^{2}(#1, #2)}
\newcommand{\Ltwonorm}[1]{\abs{\abs{#1}}_2}
\newcommand{\Lonespace}[2]{\mathscr{L}^{1}(#1, #2)}
\newcommand{\integral}[4]{\int_{#1}^{#2} #4 \, d#3} 

\section*{Q9} 
Fix some $x \in (a, b)$ and take an arbitrary sequence $(x_n)$ converging to $x$. 
Clearly $f(t) \chi_{[a, x_n]}(t) \rightarrow f(t)\chi_{[a, x]}(t)$ for every $t \in [a, b] \setminus \{x\}$. 
Further, $\abs{f_n} \leq \abs{f}$ on $[a, b]$, hence the dominated convergence theorem gives 
\begin{align*}
  F(x) = \int_{a}^{b} f\chi_{[a, x]}dt 
  & = \lim_{n \rightarrow \infty} \int_{a}^{b} f \chi_{[a, x_n]} dt \\  
  & = \lim_{n \rightarrow \infty} \int_{a}^{x_n} f dt \\ 					
  & = \lim_{n \rightarrow \infty} F(x_n)
\end{align*}
$(x_n)$ is an arbitrary sequence converging to $x$, so $F$ is continuous at $x$. 

\section*{Q10} 
Recall $f \in \Lonespace{X}{\mu} \iff \abs{f} \in \Lonespace{X}{\mu} \iff \int_{X} \abs{f} \, d\mu < \infty$. 
Let $A = \{x \in X \mid \abs{f(x)} < 1\}$ and $B = \{x \in X \mid \abs{f(x)} \geq 1\}$ partition $X$. 
\begin{align}
  \integral{X}{}{\mu}{\abs{f}} & = \integral{A}{}{\mu}{\abs{f}} + \integral{B}{}{\mu}{\abs{f}} \nonumber \\ 
							   & \leq \integral{A}{}{\mu}{1} + \integral{B}{}{\mu}{\abs{f}^2} \nonumber \\
							   & = m(A) + \Ltwonorm{f\chi_{B}}^2 \label{eq:10_sum}
\end{align}
The first term in \eqref{eq:10_sum} is finite because $A \subset E$ and $m(E) < \infty$ by assumption. 
The second term in \eqref{eq:10_sum} is also finite because $f \in \Ltwospace{X}{\mu}$ and $\abs{f\chi_B} \leq \abs{f}$ on $E$. 
This shows $f \in \Lonespace{X}{\mu}$. 

Next, let $f(x) = \frac{1}{1 + \abs{x}}$. 
Notice $f$ is symmetric about $x = 0$. 
$f$ is nonnegative everywhere and decreases monotonically on $[1, \infty)$, so one can apply the integral test to $f^2$: 
\begin{equation} \label{eq:10_integraltest} 
  \sum_{n = 1}^{\infty} f(n)^2 \text{ converges} \iff \integral{1}{\infty}{x}{f^2} \text{ converges}
\end{equation}
$\sum_{n = 1}^{\infty} f(n)^2 = \sum_{n = 1}^{\infty} \frac{1}{(1 + n)^2}$ converges, so the integral in \eqref{eq:10_integraltest} converges and is finite.  
By symmetry, $\integral{-\infty}{-1}{x}{f^2}$ converges and is finite. 
Of course, $\int_{-1}^{1} f^2 \, dx$ is finite, hence 
$$\int_{\R} \abs{f}^2 \, dx < \infty,$$

Meanwhile, a similar analysis via the integral test applied to $f$ shows 
$$\int_{\R} f \, dx = \infty$$
because $\sum_{n = 1}^{\infty} f(n) = \sum_{n = 1}^{\infty} \frac{1}{1 + n}$ diverges. 
The conclusion is that $f \in \Ltwospace{\R}{m}$ but $f \notin \Lonespace{\R}{m}$. 

\section*{Q12} 
Yes, $g$ is continuous on $[0, 1]$. 
To emphasize which variable is fixed, let 
\begin{align*}
  f^{x}(t) & = f(x, t) \text{ be a function in } y \\ 
  f^{y}(s) & = f(s, y) \text{ be a function in } x.
\end{align*}
Pick an arbitrary $x \in [0, 1]$ and take any sequence $(x_n)$ converging to $x$. 
Since $f^{y}(s)$ is continuous for any fixed $s$, 
$$f^{x_n}(t) \rightarrow f^{x}(t) \text{ upon fixing any } t \in [0, 1].$$
By assumption, $\abs{f^{x_n}(t)} \leq 1$ for all $t \in [0, 1]$, so dominated convergence theorem gives 
\begin{align*}
  g(x) = \integral{0}{1}{y}{f^{x}(y)}  
  & = \lim_{n \rightarrow \infty} \integral{0}{1}{y}{f^{x_n}(y)} \\ 
  & = \lim_{n \rightarrow \infty} g(x_n)
\end{align*}
$(x_n)$ is an arbitrary sequence converging to $x$, so $g$ is continuous at $x$. 

\section*{Q16} 
Note $E$ is measurable by Question 3 of Rudin. 
Take any measurable $A \subset E$. 
Clearly, $\chi_A \in \Ltwospace{\piInt}{m}$ and, with respect to $(e^{inx})$, has Fourier coefficients 
\begin{align*}
  c_n & \defeq \frac{1}{2\pi} \int_{-\pi}^{\pi} \chi_A e^{-inx} dx \\ 
	  & = \frac{1}{2\pi} \int_{A} (\cos{nx} - i\sin{nx}) dx 
\end{align*}
$\sum_{-\infty}^{\infty} \abs{c_n}^{2}$ converges, so $(c_n)$ and all its subsequences converge to 0. 
Namely, $c_{n_k}, c_{2n_k} \rightarrow 0$ as $k \rightarrow \infty$. 
This gives 
$$\int_{A} \sin{n_k x} dx \rightarrow 0$$
and 
\begin{align*}
  2 \int_{A} (\sin{n_k x})^2 dx & \overset{\text{trig}}{=} \int_{A} (1 - \cos{2n_k x})dx \\ 
								  & \rightarrow \int_{A} 1 dx = m(A)
\end{align*}
as $k \rightarrow \infty$, proving the two hints given by Rudin. 

Define 
$$f(x) \defeq \lim_{k \rightarrow \infty} \sin{n_k x}, \quad (x \in E).$$
Take any measurable $A \subset E$. 
On $A$, the sequence $(\sin{n_k x})_k$ still converges to $f(x)$ pointwise; 
further, the sequence is dominated by 1, so the results from the hint give 
\begin{gather}
  \int_{A} f dm = 0 \label{eq:16_hint1} \\ 
  \int_{A} f^2 dm = \frac{1}{2}m(A) \label{eq:16_hint2} 
\end{gather}
\eqref{eq:16_hint2} can be rearranged as $\int_{A} (f^2 - \frac{1}{2}) dm = 0$; 
Letting $A = E$, the result from Question 2 of Rudin says 
\begin{equation} \label{eq:16_Eae} 
  f = \pm\sqrt{\frac{1}{2}} \text{ almost everywhere on } E
\end{equation}
Let $A_1 = \{x \in E \mid f(x) = +\sqrt{\frac{1}{2}}\}$ and $A_2 = \{x \in E \mid f(x) = -\sqrt{\frac{1}{2}}\}$. 
Then 
$$\int_{A_1} f dm = \sqrt{\frac{1}{2}} m(A_1), \quad \int_{A_2} f dm = -\sqrt{\frac{1}{2}} m(A_2),$$
and considering the cases $A = A_1$ and $A = A_2$ gives $m(A_1) = m(A_2) = 0$. 
Combining this result with \eqref{eq:16_Eae} gives $m(E) = 0$, as was to be shown. 

\section*{Q17} 
Once again, clearly $\chi_E \in \Ltwospace{\piInt}{m}$. 
With respect to the orthogonal system $(e^{inx})$, it has Fourier coefficients 
\begin{align*}
  c_n & \defeq \frac{1}{2\pi} \int_{-\pi}^{\pi} \chi_{E} e^{-inx} dx \\ 
	  & = \frac{1}{2\pi} \int_{E} e^{-inx}dx \\ 
	  & = \frac{1}{2\pi} \int_{E} (\cos{nx} - \sin{nx})dx
\end{align*}

For contradiction, suppose there exists $\delta > 0$ such that infinitely many integers $n$ satisfy 
$$\forall x \in E, \quad \sin nx > \delta$$
Then for such integers $n$, 
\begin{align*}
  4\pi^2 \abs{c_n}^2 & = \left(\int_{E} \cos{nx}\,dx\right)^2 + \left(\int_{E} \sin{nx}\,dx\right)^2 \\ 
					 & \geq \delta^2 m(E)^2
\end{align*}
That there are infinitely many such $n$ contradicts a consequence of Bessel's inequality, $c_n \rightarrow 0$. 

\end{document}
